% -------------------------- 矢量分析笔记（子文件） --------------------------
% 说明：本文件为子文件，由主文档 \input{xxx} 引入，不含 document 环境与重复宏包
% 依赖宏包（请在主文档 preamble 中添加）：
% \usepackage{tcolorbox}
% \tcbuselibrary{skins}
% 定义彩色公式框：浅蓝背景、无边框、内边距适中
\newtcolorbox{formulaBox}{
    colback=blue!8!white,   % 浅蓝底色
    colframe=transparent,   % 隐藏边框
    boxrule=0pt,
    halign=center,          % 公式居中
    boxsep=2pt,
    left=4pt,right=4pt,top=3pt,bottom=3pt
}

% -------------------------- 正文 --------------------------
\section{矢量分析}

\subsection{正交坐标系}
\subsubsection{正交坐标系定义}
正交坐标系是指坐标系中各坐标基矢量相互垂直（点积为0）、且为单位矢量的坐标系，满足正交性和归一性，是矢量分析的基础框架。

\subsubsection{三种常用正交坐标系}
\begin{itemize}
    \item \textbf{直角坐标系（笛卡尔坐标系）}：坐标 \((x, y, z)\)，单位矢量 \(\vec{e}_x, \vec{e}_y, \vec{e}_z\)（常矢量），适用于规则几何形状、线性问题。
    \item \textbf{柱坐标系}：坐标 \((\rho, \phi, z)\)，单位矢量 \(\vec{e}_\rho, \vec{e}_\phi, \vec{e}_z\)（\(\vec{e}_\rho, \vec{e}_\phi\) 为变矢量），适用于圆柱对称问题。
    \item \textbf{球坐标系}：坐标 \((r, \theta, \phi)\)，单位矢量 \(\vec{e}_r, \vec{e}_\theta, \vec{e}_\phi\)（均为变矢量），适用于球对称问题。
\end{itemize}

\subsubsection{坐标变换公式}
\begin{itemize}
    \item 直角坐标 \(\leftrightarrow\) 柱坐标
    \[
    \rho = \sqrt{x^2+y^2},\quad \phi = \arctan\frac{y}{x},\quad z=z
    \]
    \item 直角坐标 \(\leftrightarrow\) 球坐标
    \[
    r = \sqrt{x^2+y^2+z^2},\quad \theta = \arccos\frac{z}{r},\quad \phi = \arctan\frac{y}{x}
    \]
\end{itemize}

\vspace{0.8cm}\hrule\vspace{0.8cm}

\subsection{矢量代数}
\subsubsection{矢量基本表示}
矢量为既有大小又有方向的量，直角坐标系中表示为：
\[
\vec{A} = A_x \vec{e}_x + A_y \vec{e}_y + A_z \vec{e}_z
\]
矢量模长：
\[
|\vec{A}| = \sqrt{A_x^2 + A_y^2 + A_z^2}
\]

\subsubsection{矢量点乘（数量积）}
\begin{itemize}
    \item \textbf{定义}：$\vec{A} \cdot \vec{B} = |\vec{A}||\vec{B}|\cos\theta$，$\theta$ 为两矢量夹角。
    \item \textbf{直角坐标（核心公式）}：
    \begin{formulaBox}
    \[
    \vec{A} \cdot \vec{B} = A_xB_x + A_yB_y + A_zB_z
    \]
    \end{formulaBox}
    \item \textbf{性质}：满足交换律、分配律，运算结果为标量。
    \item \textbf{几何意义}：表征两矢量同向程度，正交矢量点积为 $0$。
\end{itemize}

\subsubsection{矢量叉乘（矢量积）}
\begin{itemize}
    \item \textbf{定义}：$\vec{A} \times \vec{B} = |\vec{A}||\vec{B}|\sin\theta \, \vec{e}_n$，运算结果为矢量。
    \item \textbf{直角坐标（行列式，核心公式）}
    \begin{formulaBox}
    \[
    \vec{A} \times \vec{B} = \begin{vmatrix}
    \vec{e}_x & \vec{e}_y & \vec{e}_z \\
    A_x & A_y & A_z \\
    B_x & B_y & B_z
    \end{vmatrix}
    \]
    \end{formulaBox}
    \item \textbf{性质}：反交换律 $\vec{A} \times \vec{B} = -\vec{B} \times \vec{A}$，不满足结合律。
    \item \textbf{几何意义}：模长等于两矢量张成的平行四边形面积。
\end{itemize}

\vspace{0.8cm}\hrule\vspace{0.8cm}

\subsection{标量场与矢量场}
\begin{itemize}
    \item \textbf{标量场}：空间每一点对应一个标量数值，如电势 $\varphi(\vec{r})$、温度 $T(\vec{r})$。
    \item \textbf{等值面（等势面）}：满足 $\varphi(x,y,z) = C$ 的空间曲面，等值面互不相交。
    \item \textbf{矢量场}：空间每一点对应一个矢量，如电场 $\vec{E}(\vec{r})$、磁场 $\vec{H}(\vec{r})$。
    \item \textbf{矢量线方程（核心公式）}：描述矢量场方向的曲线，满足微分关系
    \begin{formulaBox}
    \[
    \frac{dx}{A_x} = \frac{dy}{A_y} = \frac{dz}{A_z}
    \]
    \end{formulaBox}
\end{itemize}

\vspace{0.8cm}\hrule\vspace{0.8cm}

\subsection{标量场的梯度}
\subsubsection{哈密顿算子}
矢量微分算子，直角坐标系形式：
\begin{formulaBox}
\[
\nabla = \frac{\partial}{\partial x}\vec{e}_x + \frac{\partial}{\partial y}\vec{e}_y + \frac{\partial}{\partial z}\vec{e}_z
\]
\end{formulaBox}

\subsubsection{梯度定义与意义}
\begin{itemize}
    \item \textbf{梯度公式（核心）}
    \begin{formulaBox}
    \[
    \nabla \varphi = \frac{\partial \varphi}{\partial x}\vec{e}_x + \frac{\partial \varphi}{\partial y}\vec{e}_y + \frac{\partial \varphi}{\partial z}\vec{e}_z
    \]
    \end{formulaBox}
    \item \textbf{物理意义}：标量场变化率最大的方向，大小等于最大变化率。
    \item \textbf{几何意义}：梯度方向与等值面垂直。
\end{itemize}

\subsubsection{方向导数}
方向导数描述标量场 $\varphi$ 在空间某一点，沿\textbf{指定方向}的变化率。
\begin{itemize}
    \item $l$：任意给定的\textbf{方向矢量}（大小为该方向的步长）。
    \item $\vec{l}_0$：$l$ 方向上的\textbf{单位矢量}，即 $\vec{l}_0 = \dfrac{\vec{l}}{|\vec{l}|}$，模长 $|\vec{l}_0|=1$。
    \item \textbf{计算公式（核心）}
    \begin{formulaBox}
    \[
    \frac{\partial \varphi}{\partial l} = \nabla \varphi \cdot \vec{l}_0
    \]
    \end{formulaBox}
    \item 最大值为 $|\nabla\varphi|$（沿梯度方向），最小值为 $-|\nabla\varphi|$（沿梯度反方向）。
    \item 沿等值面切线方向，方向导数恒为 $0$。
\end{itemize}

\vspace{0.8cm}\hrule\vspace{0.8cm}

\subsection{矢量场的通量与散度}
\subsubsection{通量}
矢量场穿过曲面的流量称为通量：
\[
\Phi = \iint_{S} \vec{A} \cdot \mathrm{d}\vec{S}
\]
闭合曲面通量（外法向为正，常用）：
\begin{formulaBox}
\[
\Phi = \oiint_{S} \vec{A} \cdot \mathrm{d}\vec{S}
\]
\end{formulaBox}

\subsubsection{散度}
散度表征空间微元内场源的发散强度，定义为单位体积闭合曲面的净通量。
\begin{itemize}
    \item \textbf{定义式}
    \[
    \nabla \cdot \vec{A} = \lim_{\Delta V \to 0} \frac{1}{\Delta V}\oiint_{S} \vec{A} \cdot \mathrm{d}\vec{S}
    \]
    \item \textbf{直角坐标（核心公式）}
    \begin{formulaBox}
    \[
    \nabla \cdot \vec{A} = \frac{\partial A_x}{\partial x} + \frac{\partial A_y}{\partial y} + \frac{\partial A_z}{\partial z}
    \]
    \end{formulaBox}
\end{itemize}

\subsubsection{柱坐标散度（重点公式）}
\begin{formulaBox}
\[
\nabla \cdot \vec{A} = \frac{1}{\rho} \frac{\partial (\rho A_\rho)}{\partial \rho} + \frac{1}{\rho} \frac{\partial A_\phi}{\partial \phi} + \frac{\partial A_z}{\partial z}
\]
\end{formulaBox}

\subsubsection{球坐标散度（重点公式）}
\begin{formulaBox}
\[
\nabla \cdot \vec{A} = \frac{1}{r^2} \frac{\partial (r^2 A_r)}{\partial r} 
+ \frac{1}{r \sin\theta} \frac{\partial (A_\theta \sin\theta)}{\partial \theta} 
+ \frac{1}{r \sin\theta} \frac{\partial A_\phi}{\partial \phi}
\]
\end{formulaBox}

\begin{itemize}
    \item \textbf{物理意义}：描述场的源/汇强度，$\nabla\cdot\vec{A}>0$ 为源，$<0$ 为汇。
\end{itemize}

\subsubsection{散度定理（高斯定理，核心定理）}
\begin{formulaBox}
\[
\oiint_{\partial V} \vec{A} \cdot \mathrm{d}\vec{S} = \iiint_V \nabla \cdot \vec{A} \, \mathrm{d} V
\]
\end{formulaBox}
闭合曲面积分 $\iff$ 体积分。

\vspace{0.8cm}\hrule\vspace{0.8cm}

\subsection{矢量场的环量与旋度}
\subsubsection{环量}
矢量场沿闭合曲线的线积分，反映场的旋转特性：
\begin{formulaBox}
\[
\Gamma = \oint_L \vec{A} \cdot \mathrm{d}\vec{l}
\]
\end{formulaBox}

\subsubsection{旋度}
\begin{itemize}
    \item \textbf{直角坐标（行列式，核心）}
    \begin{formulaBox}
    \[
    \nabla \times \vec{A} = \begin{vmatrix}
    \vec{e}_x & \vec{e}_y & \vec{e}_z \\
    \partial_x & \partial_y & \partial_z \\
    A_x & A_y & A_z
    \end{vmatrix}
    \]
    \end{formulaBox}
\end{itemize}

\begin{itemize}
    \item \textbf{柱坐标（行列式，重点）}
    \begin{formulaBox}
    \[
    \nabla \times \vec{A} = \frac{1}{\rho}
    \begin{vmatrix}
    \vec{e}_\rho & \rho\vec{e}_\phi & \vec{e}_z \\
    \frac{\partial}{\partial \rho} & \frac{\partial}{\partial \phi} & \frac{\partial}{\partial z} \\
    A_\rho & \rho A_\phi & A_z
    \end{vmatrix}
    \]
    \end{formulaBox}

    \item \textbf{球坐标（行列式，重点）}
    \begin{formulaBox}
    \[
    \nabla \times \vec{A} = \frac{1}{r^2\sin\theta}
    \begin{vmatrix}
    \vec{e}_r & r\vec{e}_\theta & r\sin\theta\vec{e}_\phi \\
    \frac{\partial}{\partial r} & \frac{\partial}{\partial \theta} & \frac{\partial}{\partial \phi} \\
    A_r & r A_\theta & r\sin\theta A_\phi
    \end{vmatrix}
    \]
    \end{formulaBox}
\end{itemize}

\subsubsection{斯托克斯定理（核心定理）}
\begin{formulaBox}
\[
\oint_L \vec{A} \cdot \mathrm{d}\vec{l} = \iint_S (\nabla \times \vec{A}) \cdot \mathrm{d}\vec{S}
\]
\end{formulaBox}

\vspace{0.8cm}\hrule\vspace{0.8cm}

\subsection{矢量恒等式与亥姆霍兹定理}
\subsubsection{重要矢量恒等式（必考公式）}
\begin{itemize}
    \item 梯度的旋度恒为零：
    \begin{formulaBox}
    \[
    \nabla \times (\nabla \varphi) = 0
    \]
    \end{formulaBox}
    \item 旋度的散度恒为零：
    \begin{formulaBox}
    \[
    \nabla \cdot (\nabla \times \vec{A}) = 0
    \]
    \end{formulaBox}
\end{itemize}

\subsubsection{亥姆霍兹定理}
\begin{itemize}
    \item 若矢量场的散度、旋度及边界条件给定，且场在无穷远处趋于0，则场唯一确定。
    \item 任意矢量场可唯一分解为：无旋场（梯度场）+ 无散场（旋度场）。
    \item 是电磁场理论与流体力学中场的唯一性与分解的基础。
\end{itemize}